\documentclass[UTF8,a4paper,10pt]{ctexart}
\usepackage[a4paper,margin=1.75cm]{geometry}
\usepackage{amsmath,amssymb,booktabs,tabularx,tikz}
\usetikzlibrary{arrows.meta,positioning}
\usepackage[most]{tcolorbox}
\setlength{\parindent}{2em}\setlength{\parskip}{0.3em}
\newtcolorbox{corebox}[1]{colback=gray!8,colframe=black!70,title=\textbf{#1},boxrule=.7pt,arc=1mm}
\begin{document}
\section*{B01\quad 局部线性化：微分 $\times$ 线性映射}
\textbf{定位：}本 Bridge 解释多元可微为何产生唯一的最佳局部线性映射，并把它交给矩阵表示。
\begin{tcolorbox}[title=母接口]
$F(x+h)=F(x)+DF_x(h)+o(\|h\|)$；可微的核心不是“偏导都存在”，而是存在一个统一线性主部。
\end{tcolorbox}
\begin{center}
\begin{tikzpicture}[node distance=9mm,every node/.style={draw,rounded corners,align=center,minimum width=3.0cm,minimum height=8mm,font=\small},>=Latex]
\node (a) {非线性映射\\$F(x+h)$};
\node (b) [right=of a] {线性主部\\$DF_x(h)$};
\node (c) [right=of b] {矩阵表示\\$J_F(x)h$};
\draw[->,thick] (a)--(b); \draw[->,thick] (b)--(c);
\end{tikzpicture}
\end{center}
\subsection*{1.\ 统一误差条件}
对 $F:\mathbb R^n\to\mathbb R^m$，可微意味着存在线性映射 $L$ 使
\[
\lim_{\|h\|\to0}\frac{\|F(x+h)-F(x)-Lh\|}{\|h\|}=0.
\]
若偏导在邻域连续，$L$ 的矩阵就是 Jacobian $J_F(x)$。链式法则变成
\[
D(G\circ F)_x=DG_{F(x)}\circ DF_x,\qquad
J_{G\circ F}=J_G(F(x))J_F(x).
\]
\subsection*{2.\ 为什么偏导存在不够}
偏导只沿坐标轴观察，不能控制所有方向的统一余项。若某点的路径 $h=t v(t)$ 使
\[
\frac{\|F(x+h)-F(x)-Lh\|}{\|h\|}
\not\to0,
\]
则不存在同一个最佳局部线性映射，即使每个坐标偏导都存在。
\subsection*{3.\ 与下游接口}
B02 取 $J_F$ 的 determinant 解释局部体积缩放；B03 取二阶余项形成 Hessian 二次型；B04 把一阶线性主部限制到切空间。矩阵只是表示，不能把 $F$ 本身误认成线性。
\subsection*{4.\ 最小例子与调用}
令 $F(x,y)=(x^2+xy,e^y)$。在 $(1,0)$ 处
\[
J_F(1,0)=\begin{pmatrix}2&1\\0&1\end{pmatrix},\qquad
F(1+h,k)=F(1,0)+J_F(1,0)\binom hk+o(\sqrt{h^2+k^2}).
\]
题面要求“小增量近似、误差传播、复合映射导数”时，先寻找这个统一线性主部，而不是逐坐标猜变化。
\begin{corebox}{检查句}
我是否控制了所有方向的余项？矩阵的行列尺寸是否等于输入/输出维度？复合时因子顺序是否保持？
\end{corebox}
\end{document}
